\documentclass[12pt,a4paper]{article}

% ============================================================
% Пакеты
% ============================================================
\usepackage[utf8]{inputenc}
\usepackage[T2A]{fontenc}
\usepackage[russian]{babel}
\usepackage{amsmath,amssymb}
\usepackage{geometry}
\usepackage{enumitem}
\usepackage{booktabs}

\geometry{left=2.5cm, right=2.5cm, top=2.5cm, bottom=2.5cm}

\pagestyle{empty}

\begin{document}

% ============================================================
% Тема
% ============================================================
\textbf{\large Тема}

\medskip
Улучшенная нижняя оценка для константы Чватала--Санкоффа методом динамического программирования с окнами

\bigskip

% ============================================================
% Краткое описание предметной области
% ============================================================
\textbf{\large Краткое описание предметной области}

\medskip
Наибольшая общая подпоследовательность (Longest Common Subsequence, LCS) двух строк~--- одна из фундаментальных задач комбинаторики и теоретической информатики, лежащая в основе алгоритмов выравнивания биологических последовательностей, утилит сравнения файлов (\texttt{diff}) и систем контроля версий.

В 1975 году Чватал и Санкофф показали, что для двух независимых случайных бинарных строк длины~$n$ (каждый бит выбирается равновероятно) математическое ожидание длины LCS растёт линейно:
\[
  \mathbb{E}[L_n] \;\sim\; \gamma_2 \cdot n, \qquad n \to \infty,
\]
где $\gamma_2$~--- \emph{константа Чватала--Санкоффа} для бинарного алфавита. Точное значение этой константы остаётся неизвестным уже 50 лет. Задача её определения тесно связана с задачей Улама об увеличивающихся подпоследовательностях случайных перестановок и с теорией случайных матриц.

Сложность задачи обусловлена тем, что замкнутая формула для $\mathbb{E}[L_n]$ не существует даже при малых~$n$, а комбинаторные методы быстро упираются в экспоненциальный рост числа конфигураций. Прогресс в оценках~$\gamma_2$ за последние десятилетия был медленным: лучшая нижняя оценка составляла $\gamma_2 \geq 0{,}792666$ (Heineman et al., 2024), а лучшая верхняя~--- $\gamma_2 \leq 0{,}826280$ (Lueker, 2009). Разрыв между оценками~$\approx 0{,}034$ делает задачу сужения этого интервала актуальной как с теоретической, так и с вычислительной точки зрения.

\bigskip

% ============================================================
% Текущее состояние и проблемы
% ============================================================
\textbf{\large Текущее состояние и проблемы}

\begin{itemize}[leftmargin=2em]
  \item Полный перебор всех $4^n$ пар строк для вычисления $\mathbb{E}[L_n]/n$ работает за $O(4^n \cdot n^2)$, что делает его неприменимым при~$n > 14{-}20$ на современном оборудовании.
  \item Существующие методы нижних оценок (конечные автоматы, марковские стратегии) достигли потолка при~$\gamma_2 \geq 0{,}7927$ и требуют экспоненциально растущих вычислительных ресурсов для каждого дальнейшего улучшения.
  \item Эмпирические оценки методом Монте-Карло указывают на значение $\gamma_2 \approx 0{,}811$, что свидетельствует о значительном пространстве для улучшения нижней оценки.
  \item Масштабирование существующих методов затруднено экспоненциальным ростом числа состояний, что требует разработки специализированных алгоритмических и программных оптимизаций.
\end{itemize}

\bigskip

% ============================================================
% Цель
% ============================================================
\textbf{\large Цель}

\medskip
Получить улучшенную строгую нижнюю оценку для константы Чватала--Санкоффа~$\gamma_2$, превосходящую все опубликованные результаты, с помощью нового вычислительного метода на основе динамического программирования с ограниченным окном обзора.

\bigskip

% ============================================================
% Задачи
% ============================================================
\textbf{\large Задачи}

\begin{itemize}[leftmargin=2em]
  \item Разработать вычислительный метод оценки $\mathbb{E}[\mathrm{LCS}]$ на основе динамического программирования с ограниченным окном обзора: на каждом шаге стратегия <<видит>> следующие~$N$ символов обеих строк и выбирает действие, максимизирующее ожидаемую длину общей подпоследовательности.
  \item Реализовать данный метод на языке C++ с оптимизациями, обеспечивающими работу с большими параметрами: итеративный подход вместо рекурсии, использование \texttt{float} вместо \texttt{double}, компиляторные оптимизации (AVX2, \texttt{-Ofast}, развёртка циклов).
  \item Провести масштабные численные эксперименты для получения нижних оценок при различных значениях параметров ($N$ от~5 до~12, $\mathit{MAX\_SH}$ до~30, $K$ до~$40\,000$).
  \item Выполнить независимую верификацию полученных результатов методом Монте-Карло на строках длиной до~$20\,000$ символов.
  \item Сравнить полученные результаты с лучшими известными оценками из литературы (Чватал--Санкофф 1975, Данчик--Патерсон 1995, Люкер 2003, Heineman et al. 2024).
\end{itemize}

\bigskip

% ============================================================
% Полученные результаты
% ============================================================
\textbf{\large Полученные результаты}

\begin{itemize}[leftmargin=2em]
  \item Получена новая лучшая нижняя оценка: $\gamma_2 \geq 0{,}7964$ при параметрах $N = 12$, $\mathit{MAX\_SH} = 30$, $K = 40\,000$. Это превосходит предыдущий рекорд $0{,}7927$ (Heineman et al., 2024) на $\approx 0{,}004$, что сопоставимо с суммарным прогрессом в данной области за период 2003--2024 гг.
  \item Реализация на C++ прошла несколько итераций оптимизации, что позволило проводить масштабные вычисления на одной машине.
  \item Дополнительно проведена эмпирическая оценка методом Монте-Карло (на строках длиной до $20\,000$), давшая значение $\gamma_2 \approx 0{,}811$, что согласуется с полученной нижней оценкой и подтверждает непротиворечивость результатов.
\end{itemize}

\medskip

\begin{table}[h]
\centering
\begin{tabular}{llc}
  \toprule
  \textbf{Авторы} & \textbf{Год} & \textbf{Нижняя оценка $\gamma_2$} \\
  \midrule
  Chv\'{a}tal, Sankoff & 1975 & $\geq 0{,}7615$ \\
  Dan\v{c}\'{i}k, Paterson & 1995 & $\geq 0{,}7739$ \\
  Lueker & 2003 & $\geq 0{,}7881$ \\
  Heineman et al. & 2024 & $\geq 0{,}7927$ \\
  \textbf{Данная работа} & \textbf{2025} & $\geq \mathbf{0{,}7964}$ \\
  \bottomrule
\end{tabular}
\end{table}

\end{document}
